\begingroup
\chapter*{Sujet de synthèse B}\markboth{Sujet de synthèse B}{Sujet de synthèse B}\addcontentsline{toc}{chapter}{Sujet de synthèse B}
\renewcommand{\thebmtexo}{SYN2.\arabic{bmtexo}}
\renewcommand{\theHbmtexo}{SYN2.\arabic{bmtexo}}
\setcounter{bmtexo}{0}
\begin{devoir}[3 h — 20 points]
Calculatrice autorisée. Les repères utilisés pour les calculs métriques sont orthonormés; ils sont directs pour le produit vectoriel. Justifier les réponses et préciser les domaines utiles. Les questions sont indépendantes sauf indication. 
\end{devoir}
\exo[Algèbre et numération]\label{t11s-syn2-e01}\bareme{6 points}
\begin{enumerate}\item Résoudre $x^3-2x^2-x+2=0$ en utilisant une racine vérifiée (2 points).
\item Calculer PGCD et PPCM de $180$ et $252$ (2 points).
\item Convertir $(110101)_2$ en base dix puis en base cinq (2 points).\end{enumerate}
\begin{corrige}[exercice~\ref{t11s-syn2-e01}]\label{sol-t11s-syn2-e01}
1. $1$ est racine; factorisation $(x-1)(x-2)(x+1)$, solutions $-1,1,2$.
2. Décompositions $2^2 3^2 5$ et $2^2 3^2 7$, PGCD $36$, PPCM $1260$.
3. Valeur $53$, soit $(203)_5$. Deux points par sous-question.
\end{corrige}
\exo[Trigonométrie et espace]\label{t11s-syn2-e02}\bareme{8 points}
\begin{enumerate}\item Résoudre $\cos x=\sqrt2/2$ sur $[-\pi;\pi]$ (2 points).
\item Calculer $\vecu\wedge\vecv$ pour $\vecu=(1;0;1),\vecv=(0;2;1)$ et l’aire du triangle engendré (2 points).
\item Déterminer le plan passant par l’origine et dirigé par ces vecteurs (2 points).
\item Trouver son intersection avec la droite $(x;y;z)=(1;0;t)$ (2 points).\end{enumerate}
\begin{corrige}[exercice~\ref{t11s-syn2-e02}]\label{sol-t11s-syn2-e02}
1. $x=\pm\pi/4$.
2. Vectoriel $(-2;-1;2)$, norme $3$, aire $3/2$.
3. Normal non nul, équation $-2x-y+2z=0$.
4. $-2+2t=0$, $t=1$, point $(1;0;1)$, unique. Deux points par sous-question.
\end{corrige}
\exo[Données et tirage]\label{t11s-syn2-e03}\bareme{6 points}
\begin{enumerate}\item Pour les valeurs $5,10,15$ d’effectifs $2,5,3$, calculer moyenne et écart type (2 points).
\item Un prix augmente de $10\%$ deux fois. Calculer le taux global (1 point).
\item Tirer uniformément trois boules simultanément parmi quatre rouges et trois bleues; calculer la probabilité d’exactement deux rouges (2 points).
\item Donner celle d’au moins une rouge (1 point).\end{enumerate}
\begin{corrige}[exercice~\ref{t11s-syn2-e03}]\label{sol-t11s-syn2-e03}
1. Moyenne $10{,}5$, variance $12{,}25$, écart type $3{,}5$.
2. Coefficient $1{,}1^2=1{,}21$, hausse $21\%$.
3. Total $35$, favorables $18$, probabilité $18/35$.
4. Contraire : trois bleues, une combinaison; donc $1-1/35=34/35$.
\end{corrige}
\endgroup
